%%%%

%%%   Самарский государственный технический университет

%%%   Кафедра прикладной математики и информатики

%%%

%%%   Преамбула для статьи в журнал "Вестник Самарского государственного

%%%   технического университета. Серия: «Физико-математические науки»"

%%%   Ориентирована под MikTEx 2.4 for Win32

%%%

%%%   Хотя сам журнал собирается под GNU/Linux на дистрибутиве TeXLive



\documentclass[11pt,a4paper,twoside]{article}



\usepackage[cp1251]{inputenc}

\usepackage[T2A]{fontenc}

\usepackage[english,russian]{babel}

\usepackage{samgtubib}

\usepackage[dvips]{graphicx}

\usepackage[multiple]{footmisc}



\usepackage{samgtu}

\renewcommand{\VestnikYear}{2009} % Год издания Вестника

\renewcommand{\VestnikNumber}{2\,(19)} % Номер Вестника

\renewcommand{\VestnikMonth}{Ноябрь} % Месяц выхода Вестника

\renewcommand{\VestnikData}{01.11.09} % Дата подписания Вестника в печать

\renewcommand{\VestnikUPL}{26,64} % Усл. печ. л.

\renewcommand{\VestnikUIL}{25,73} % Уч.-изд. л.

\renewcommand{\VestnikRegNumber}{479/09} % Регистрационный номер

\renewcommand{\VestnikOrder}{994} % Номер заказа







%%\samgtupubl % для генерации pdf, отдаваемого в типографию СамГТУ

%\pdfcrop % обрезка pdf для размещения на math-net.ru

%\draft % На корректуре

\nofilllastpage % Последняя страница не заполнена не полностью

%\briefcomm % Статья является кратким сообщением







\vsgtu{1202} % Номер статьи по базе Math-Net.ru



\FirstPage{86}

\LastPage{92}





%\begin{document}







%%%%%%%%%%%%%%%%%%%%%%%%%%%%%%%%%%%%%%%%%%%%%%%%%%%%%%%%%%%%%%%%%%%%%%%%%%%%%%%%%%%%

%Информация о статье и об авторах на русском и английском языках



% Номер УДК

\udk{539.376}%Обработка текста. Подготовка текстов : Языки программирования. Компьютерные языки

\msc{74D10, 74S60; 74R20, 74E35, 74E05}% Коды MSC см. http://www.ams.org/msc/



\rutitle{Оценка над\"eжности стохастически неоднородной~толстостенной трубы по~критерию~длительной прочности}% Полный заголовок

{Оценка надёжности стохастически неоднородной толстостенной трубы\dots}% Заголовок, который идёт в верхний колонтитул



\ruauthor{Н.\,Н.~Попов, Л.\,В.~Коваленко}{Н.~Н.~П\,о\,п\,о\,в, Л.~В.~К\,о\,в\,а\,л\,е\,н\,к\,о}% Автор(ы) для оглавления и колонтитула через \,



\ruaffil{\rusamgtu.}





\email{2}{\mem{ponick25@gmail.com} (N.N.~Popov, \CAut), \newline  \mem{flytitmouse@mail.ru} (L.V.~Kovalenko)} %E-mail's авторов



\ruauthordetails{2}{\fullauthor{\rupid{38440}{Николай Николаевич Попов}}\regalia{к.ф.-м.н.,

доц.}\position{доцент} \dept{каф. прикладной математики и информатики} 

\fullauthor{\rupid{41434}{Людмила Викторовна Коваленко}} \regalia{к.ф.-м.н.} \position{ассистент} \dept{каф.

прикладной математики и информатики}}



\rukeywords{установившаяся ползучесть, толстостенная труба,

микронеоднородный материал, стохастическая краевая задача, длительная

прочность, функция надёжности}



\ruabstract{Разработана методика вероятностной оценки надёжности микронеоднородной толстостенной трубы на основе решения стохастической краевой задачи ползучести. 

Реологические свойства материала при этом описывались при помощи случайной функции одной

переменной (радиуса $r$). 

Для изучения процесса разрушения материала при ползучести введён параметр повреждённости $0\leq\omega(t)\leq1$ и принята степенная зависимость скорости изменения $\omega(t)$ от эквивалентного напряжения $\sigma_{\text{э}}$, которое определялось по критерию Сдобырева.

Оценка надёжности производится по интегрально"=среднему значению эквивалентного напряжения. 

Найдено случайное время до разрушения и~его функция распределения, которая аппроксимировалась

логарифмически нормальным законом. 

Приведён пример вычисления вероятности безотказной работы для толстостенной трубы из микронеоднородного материала.

Полученные результаты позволяют оценивать надёжность стохастически неоднородных осесимметричных элементов конструкций при условии, что из эксперимента будут получены  необходимые статистические данные.

}



\entitle{Reliability Evaluation of Stochastically Heterogeneous Thick-walled Pipe by~Long-term~Strength Criterion}





\enauthor{N.~N.~Popov, L.~V.~Kovalenko}{N.~N.~P\,o\,p\,o\,v, L.~V.~K\,o\,v\,a\,l\,e\,n\,k\,o}





\enaffil{\ensamgtu.}



\enauthordetails{2}{\fullauthor{\enpid{38440}{Nikolay N. Popov}} \regalia{Cand. Phys. \& Math. Sci.} \position{Associate Professor} \dept{Dept. of Applied Mathematics \& Computer Science} \fullauthor{\enpid{41434}{Ludmila V. Kovalenko}} \regalia{Cand. Phys. \& Math. Sci.} \position{Assistant} \dept{Dept. of Applied Mathematics

\& Computer Science}}





\enabstract{We have developed a method of probabilistic reliability evaluation of microheterogeneous thick-walled pipe, based on the already received solution of the stochastic creep boundary value problem. The rheological properties of the material were described using random function of one variable (radius $r$). Damage parameter $0\leq\omega(t)\leq1$ was introduced here to study the process of degradation of the material during creep stage. Also the power law of the rate of $\omega(t)$ change on the equivalent stress $\sigma_{\text{e}}$, determined by Sdobyrev criterion, is assigned. The reliability evaluation is made by the mean integral value of the equivalent stress. We have found a random time before destruction and its distribution function, which was approximated by lognormal law. The problem of the probability of failure-free operation was calculated for a thick-walled microheterogeneous pipe with the specified parameters. The obtained results allow to evaluate reliability of stochastically inhomogeneous axisymmetric structural elements if necessary statistical data are obtained from the experiment.

}

%

\enkeywords{steady-creep state, thick-walled pipe, microheterogeneous material, stochastic boundary value problem, long-term strength, reliability function}



\date{20/XII/2013} % Дата поступления статьи в редакцию.

\revisiondate{21/II/2014} % В окончательном виде

\accepteddate{26/II/2014}



%%%%%%%%%%%%%%%%%%%%%%%%%%%%%%%%%%%%%%%%%%%%%%%%%%%%%%%%%%%%%%%%%%%%%%%%%%%%%%%%%%%%





\makerutitle



%Каждый пункт статьи должен начинаться с команды Название пункта

%после названия точку ставить не нужно. Пункты автоматически нумеруются.

%Если пункт нумеровать не нужно, то следует использовать в команде

% \Section необязательный параметр [n]:

%   \Section[n]{Имя пункта}

% Введение и заключение обычно не нумеруются.





В условиях длительного действия нагрузки при повышенной температуре с

некоторого момента скорость ползучести начинает возрастать (третья

стадия), и~процесс ползучести заканчивается разрушением. В силу

стохастической неоднородности материала время до разрушения будет

являться случайной величиной. Для описания процесса разрушения при

ползучести введем параметр поврежденности материала

$0\leq\omega(t)\leq1$ и используем кинетическое уравнение

Работнова \cite{popovkoval:bib:Rabotnov}.

\begin{equation}

\frac{d\omega}{dt}=B\Bigl(\frac{\sigma_{\text{э}}}{1-\omega}\Bigr)^k,

\label{eq:kov:kineticEq}

\end{equation}

где $\sigma_{\text{э}}$ "---  эквивалентное напряжение, задаваемое в

виде некоторого соотношения между компонентами тензора напряжений;

$B$, $k$ "---  постоянные материала.



В данной статье рассматривается оценка надёжности по критерию

длительной прочности толстостенной трубы из микронеоднородного

материала, находящейся под действием внутреннего давления $q$.

Предполагается, что в трубе осуществляется плоская деформация и в

традиционной цилиндрической системе координат полагается, что

компонента тензора деформаций $\varepsilon_z=0$. При этом

реологические свойства материала описываются при помощи случайной

функции радиуса $r$.



В работе \cite{popovkoval:bib:RadchPopov} рассматривалось решение

стохастической краевой задачи установившейся ползучести

толстостенной трубы. Определяющие соотношения для деформаций

ползучести $\varepsilon_r$ и $\varepsilon_\varphi$ принимаются в

соответствии с теорией вязкого течения в стохастической форме

\[

\dot {\varepsilon }_\varphi = - \dot {\varepsilon }_r = cS^{n -

1}(\sigma _\varphi - \sigma _r )(1 + \alpha U),

\]

где $\sigma _r$ и $\sigma _\varphi$ "--- радиальная и

тангенциальная компоненты тензора напряжений,  $S =

({\sqrt{3}}/{2})(\sigma _\varphi - \sigma _r )$ "--- 

интенсивность напряжений, $U(r)$ "---  случайная однородная функция,

описывающая флуктуации реологических свойств материала с

математическим ожиданием $\langle U\rangle=0$ и дисперсией

$\langle U^2\rangle=1$, $\alpha$ "---  коэффициент вариации этих

свойств (малый параметр), $c$, $n$ "---  постоянные материала.



Аналогичная задача по оценке надёжности толстостенной трубы по

критерию деформационного типа, когда ограничение накладывалось на

случайное радиальное перемещение, рассматривалась в работе

\cite{popovkoval:bib:PopovKoval}.



Случайные напряжения $\sigma _r$ и $\sigma _\varphi$, найденные в

работе \cite{popovkoval:bib:RadchPopov} на основе первого приближения метода малого

параметра, имеют вид:



\begin{equation}

\sigma _{r} = A\Bigl( 1 - \Bigl( {\frac{\beta}{r}} \Bigr)^{2/n}

\Bigr)+\frac{2A{\beta}^{2/n}\alpha}{n^2}\left[

(1-r^{-2/n})H - I(r)\right],

\label{eq:kov:2}

\end{equation}

\begin{multline}

\sigma _{\varphi} = 

A\Bigl[ 1 + \Bigl(\frac{2}{n}-1\Bigr)\Bigl( {\frac{\beta}{r}} \Bigr)^{2/n} \Bigr] + \\ +\frac{2A{\beta}^{2/n}\alpha}{n^2}

\Bigl[

\Bigl(1+\Bigl(\frac{2}{n}-1\Bigr)r^{-2/n}\Bigr)H - I(r)-Ur^{-2/n}\Bigr], \label{eq:kov:3}

\end{multline}

\[

A = \frac{q}{\beta^{2/n}-1}, \quad I (r) =

 \int _1^r {U\left( x \right)x^{ - 1 -2/n}dx},  \quad H =

\frac{I ( \beta )}{1 - {\beta}^{-2/n}}.

\]

Здесь $r$ "---  безразмерный радиус $(r\in[1, \beta])$, $\beta={b}/{a}$, $a$ и

$b$ "---  внутренний и внешний радиусы трубы, $q$ "---  давление.



В работе \cite{popovkoval:bib:RadchPopov} также найдены статистические характеристики

случайных напряжений $\sigma _r$ и $\sigma _{\varphi}$ c

математическими ожиданиями

\begin{gather}

\langle\sigma

_r\rangle=A\Bigl(1-\Bigl(\frac{\beta}{r}\Bigr)^{2/n}\Bigr),

\quad

\langle\sigma _{\varphi}\rangle=A\Bigl[1+\Bigl(\frac{2}{n}-1\Bigr)\Bigl(\frac{\beta}{r}\Bigr )^{2/n}\Bigr]

\end{gather}

и дисперсиями

\begin{gather}

D[\sigma_r]=\frac{4A^2\beta^{4/n}\alpha^2}{n^4}

\bigl[(1-r^{-2/n})^2\langle

H^2\rangle-2 (1-r^{-2/n})\langle

I(r)H\rangle+\langle I^2(r)\rangle\bigr],\label{eq:kov:4}

\end{gather}

\begin{multline}

D[\sigma_{\varphi}]=\frac{4A^2\beta^{4/n}\alpha^2}{n^4}

\Bigl[\Bigl(1+\Bigl(\frac{2}{n}-1\Bigr)r^{-2/n}\Bigr)^2

\langle H^2\rangle+ \\ + \langle I^2(r)\rangle+r^{-4/n}

-2\Bigl(1+\Bigl(\frac{2}{n}-1\Bigr)r^{-2/n}\Bigr)

\langle I(r)H\rangle-

\\

-2\Bigl(1+\Bigl(\frac{2}{n}-1\Bigr)r^{-2/n}\Bigr)\langle HU\rangle r^{-2/n}+2r^{-2/n}\langle

I(r)U\rangle\Bigr],\label{eq:kov:5}

\end{multline}



\noindent

где 

\begin{gather*}

\langle

H^2\rangle=\frac{1}{\left(1-\beta^{-2/n}\right)^2} \int _1^{\beta} \int _1^{\beta}{K(x_2-x_1)x_1^{-1-2/n}x_2^{-1-2/n}dx_1dx_2},

\\

\langle

I^2(r)\rangle= \int _1^{r} \int _1^{r}{K(x_2-x_1)x_1^{-1-2/n}x_2^{-1-2/n}dx_1dx_2},\\

\langle

I(r)H\rangle=\frac{1}{1-\beta^{-2/n}} \int _1^{r} \int _1^{\beta}{K(x_2-x_1)x_1^{-1-2/n}x_2^{-1-2/n}dx_1dx_2},

\\

\langle

HU\rangle=\frac{1}{1-\beta^{-2/n}} \int _1^{\beta}{K(x-r)x^{-1-2/n}dx},\\

\langle I(r)U\rangle= \int _1^{r}{K(x-r)x^{-1-2/n}dx},

\end{gather*}

$K(x_2-x_1)$ "---  корреляционная функция случайной функции $U(r)$.



Эквивалентное напряжение $\sigma_{\text{э}}$, входящее в

кинетическое уравнение \eqref{eq:kov:kineticEq}, будем определять

по критерию Сдобырева, который довольно часто применяется для

описания закономерности изменения длительной прочности при сложном

напряженном состоянии \cite{popovkoval:bib:Sdobyrev}:

\[

\sigma_{\text{э}}=\frac{\sigma_1+S}{2},

\]

где $\sigma_1$ "---  наибольшее нормальное напряжение, которое для

толстостенной трубы равно $\sigma _\varphi$.



В общем случае эквивалентное напряжение $\sigma_{\text{э}}$ является случайной функцией радиуса $r$ и времени $t$. 

Её характеристики можно найти из соответствующей стохастической краевой задачи, при решении которой возникают непреодолимые трудности. 

В~связи с этим предлагается использовать введённое в работе  \cite{popovkoval:bib:Lokos} интегрально"=среднее значение эквивалентного напряжения:

\begin{equation}

\overline{\sigma}_{\text{э}}=\frac{1}{\beta-1} \int _1^{\beta}{\sigma_{\text{э}}(r)dr}.\label{eq:kov:9}

\end{equation}



В работе \cite{popovkoval:bib:RadBash}   показано, что в детерминированной постановке

$\sigma_{\text{э}}$ является практически постоянной для  толстостенной трубы при ползучести от момента $t=0$ вплоть до разрушения для широкого диапазона изменения параметров $n$ и $\beta$. 

Поэтому вводится гипотеза, что и  в стохастической постановке $\overline{\sigma}_{\text{э}}$ считается случайной величиной, но не зависящей от времени, которую можно вычислить, например, на стадии установившейся ползучести.  

В~дальнейшем теория длительной прочности строится на основе интегрально"=среднего значения  эквивалентного напряжения $\overline{\sigma}_{\text{э}}$, статистические характеристики которого могут быть найдены исходя из \eqref{eq:kov:2} и \eqref{eq:kov:3}.



Согласно \eqref{eq:kov:2}, \eqref{eq:kov:3} и \eqref{eq:kov:9}

величина эквивалентного напряжения $\overline{\sigma}_{\text{э}}$

определяется формулой



\begin{multline}

\overline{\sigma}_{\text{э}}=\frac{1}{2(\beta-1)} 

\int _1^{\beta}

\Bigl(\sigma_{\varphi}+\frac{\sqrt{3}}{2}(\sigma _\varphi - \sigma _r )\Bigr)dr=\\

=\frac{A}{2}

\Bigl[1 +\frac{\beta^{2/n}\Bigl(1-\frac{2}{n}\Bigr)}{\beta-1}

\Bigl(\frac{2+\sqrt{3}}{n}-1\Bigr)

\Bigl( \frac{1}{\beta^{1-2/n}} -1\Bigr)\Bigr]+\label{eq:kov:6}\\

+\frac{A{\beta}^{2/n}\alpha}{n^2(\beta-1)} 

\int _1^{\beta}\Bigl[

\Bigl(1+\Bigl(\frac{2+\sqrt{3}}{n}-1\Bigr)r^{-2/n}\Bigr)H - \\ -

I(r)-\Bigl(1+\frac{\sqrt{3}}{2}\Bigr)Ur^{-2/n}\Bigr]dr.

\end{multline}





%Оценка локальной прочности в соответствии с \cite{popovkoval:bib:RabotnovMileyko} производится

%по наибольшему значению эквивалентного напряжения

%$\sigma_{\text{э}max}$. При этом мы получим нижнюю оценку

%долговечности элемента конструкции. Распределение абсолютного

%максимума $\sigma_{\text{эmax}}$ случайной функции

%$\sigma_{\text{э}}(r)$ можно найти приближенно по теории выбросов

%\cite{popovkoval:bib:Bolotin}. Однако случайная функция $\sigma_{\text{э}}(r)$ является

%неоднородной, поэтому распределение абсолютного максимума будет

%задаваться громоздким выражением, которое затруднительно применять

%для вероятностной оценки времени до разрушения. В связи с этим

%вводится упрощающее предположение, что разрушение начинается с

%точки, в которой максимально среднее значение эквивалентного

%напряжения

%\[

%\langle\sigma_{\text{э}}\rangle=\frac{A}{2}\left[1+\left(\frac{2+\sqrt{3}}{n}-1\right)\left(\frac{\beta}{r}\right)^{2/n}\right]

%\]

%

%При $n<2+\sqrt{3}$ величина $\langle\sigma_{\text{э}}\rangle$

%достигает наибольшего значения на внутреннем радиусе трубы, откуда

%и начинается разрушение. При $n>2+\sqrt{3}$ максимум

%$\langle\sigma_{\text{э}}\rangle$ достигается на наружнем радиусе.

%Если $n=2+\sqrt{3}$, то в состоянии установившейся ползучести

%величина $\langle\sigma_{\text{э}}\rangle$ постоянна по толщине

%трубы.



Решая дифференциальное уравнение \eqref{eq:kov:kineticEq} при

$\omega(0)=0$ и считая, что в момент разрушения $\omega(t_{\rm p})=1$,

можно найти случайное время до разрушения

\[

t_{\rm p}=\frac{1}{B(k+1)\overline{\sigma}^k_{\text{э}}}.

\]



Время до разрушения $t_{\rm p}$ было аппроксимировано логарифмически

нормальным законом, функция распределения которого имеет вид

\begin{equation}

F(t)=\frac{1}{\sqrt{2\pi}d} \int _0^t{\frac{1}{z}\exp{\Bigl[-\frac{\left(\ln

z-m\right)^2}{2d^2}\Bigr]dz}}. \label{eq:kov:7}

\end{equation}



Для того чтобы найти параметры  $m$ и $d$ распределения \eqref{eq:kov:7}, разложим величину $\ln{\sigma_{\text{э}}}$ по малому

параметру $\alpha$ и ограничимся членами только первого порядка относительно $\alpha$:

\begin{gather}

\ln{\sigma_{\text{э}}}=

\ln\bigl(\overline{\sigma}_{\text{э}}^0+\alpha\overline{\sigma}_{\text{э}}^{\ast}\bigr)=

\ln{\overline{\sigma}_{\text{э}}^0}+\ln{\Bigl(1+\frac{\alpha\overline{\sigma}_{\text{э}}^{\ast}}{\overline{\sigma}_{\text{э}}^0}\Bigr)}=

\ln{\overline{\sigma}_{\text{э}}^0}+\frac{\alpha\overline{\sigma}_{\text{э}}^{\ast}}{\overline{\sigma}_{\text{э}}^0}+{\sf o}(\alpha),\label{eq:kov:8}

\end{gather}

где

$$

\overline{\sigma}_{\text{э}}^0=\langle\overline{\sigma}_{\text{э}}\rangle=

\frac{A}{2}\Bigl[1

+\frac{\beta^{2/n}\left(1-\frac{2}{n}\right)}{\beta-1}

\Bigl(\frac{2+\sqrt{3}}{n}-1\Bigr)

\Bigl( \frac{1}{\beta^{1-2/n}} -1\Bigr)\Bigr],\quad 

\langle\overline{\sigma}_{\text{э}}^{\ast}\rangle=0.$$ 

Используя разложение \eqref{eq:kov:8}, вычислим параметры распределения

\eqref{eq:kov:7}:

\begin{gather}

m=\langle\ln{t_p}\rangle=-\ln{\left(B(k+1)\right)}-k\ln{\overline{\sigma}_{\text{э}}^0},\quad 

d^2=D\left[\ln{t_p}\right]=\frac{k^2\alpha^2}{\overline{\sigma}_{\text{э}}^0}D\left[\overline{\sigma}_{\text{э}}^{\ast}\right],\notag

\end{gather}

где дисперсия $D\left[\overline{\sigma}_{\text{э}}^{\ast}\right]$ определяется 

согласно \eqref{eq:kov:6}:

$$

D\left[\overline{\sigma}_{\text{э}}^{\ast}\right]=

D\left[\overline{\sigma}_{\text{э}}\right]=

\frac{1}{\beta-1} \int _1^{\beta}

\Bigl[\frac{1}{4}\Bigl(1+\frac{\sqrt{3}}{2}\Bigr)^2D\left[\sigma_{\varphi}\right]+\frac{3}{16}D\left[\sigma_r\right]-\frac{\sqrt{3}}{2}\langle\sigma_{\varphi}\sigma_r\rangle\Bigr]dr.

$$

Здесь величины $D\left[\sigma_{\varphi}\right]$ и

$D\left[\sigma_r\right]$ вычисляются по формулам \eqref{eq:kov:4}

и \eqref{eq:kov:5} соответственно. 

Момент

$\langle\sigma_{\varphi}\sigma_r\rangle$ можно найти, используя \eqref{eq:kov:2} и \eqref{eq:kov:3}:

\begin{multline}

\langle\sigma_{\varphi}\sigma_r\rangle=

\frac{4A^2\beta^{4/n}\alpha^2}{n^2}

\Bigl[(1-r^{-2/n})

\Bigl(1+\Bigl(\frac{2}{n}-1\Bigr)r^{-2/n}\Bigr)\langle H^2\rangle-

\\

-2\Bigl(1+\Bigl(\frac{1}{n}-1\Bigr)r^{-2/n}\Bigr)\langle

I(r)H\rangle+\langle I^2(r)\rangle+\\

+\langle I(r)U\rangle r^{-2/n}-\bigl(1-r^{-2/n}\bigr)r^{-2/n}\langle HU\rangle.

\end{multline}



Рассмотрим модельную задачу оценки надёжности для толстостенной трубы из стали 12ХМФ при $T=590$\,\textcelsius\ с~постоянными материала $n=7.1$, $c=2.762\cdot10^{-21}$, внутренним и внешним радиусами

$a=14$ мм и $b=21$ мм соответственно, находящейся под действием внутреннего давления $q=70$~МПа. Исходные данные для расчёта взяты из монографии \cite{popovkoval:bib:RadchEremin}. 

Параметры кинетического уравнения \eqref{eq:kov:kineticEq} считались равными $B=2.54\cdot10^{-12}$,

$k=3.04$.

Корреляционная функция случайной функции $U(r)$ была взята в виде

$$K(r)=\exp{(-10|r|)}\left(\cos{20r}+0.5\sin{20|r|}\right), \quad  r=x_2-x_1.$$



На \hyperlink{popov:pic}{рисунке} представлена функция надёжности

$P(t)=1-F(t)$, где $F(t)$ определяется по формуле \eqref{eq:kov:7}

при степени неоднородности материала $\alpha=0.3$. Функцию

надёжности $P(t)$ можно использовать для назначения ресурса

толстостенной трубы. Назначенный ресурс $T_{\ast}$ определяют так,

чтобы вероятность обеспечения $T_{\ast}$ была  равна заданному

значению $p_{\ast}$ вероятности безотказной работы (пунктирные

линии на графике). При заданном значении $p_{\ast}=0.95$ ресурс

для рассматриваемой трубы составляет 19\,475~часов.



\begin{figure}[h!]\centering \small

\includegraphics[scale=0.9]{lu} \hypertarget{popov:pic}{}

\caption*{Функция надёжности $P(t)$ [Reliability function {\it vs} time]}

\end{figure}



Полученные результаты позволяют оценивать надёжность стохастически неоднородных осесимметричных элементов конструкций при условии, что необходимые статистические данные будут получены экспериментальным путём.







\thanks{Работа выполнена при поддержке РФФИ (проект \No~13--01--00699--a).}



\thanks{This work is supported by RFBR, project no. 13--01--00699--a.}



\pagebreak



%% Благодарности, гранты и прочее



\begin{thebibliography}{9}



\RBibitem{popovkoval:bib:Rabotnov} 

\by Ю.~Н.~Работнов 

\book Ползучесть элементов конструкций 

\publaddr М. 

\publ Наука 

\yr 1966

\totalpages 752

\misctransl

\lang In Russian

\by Yu.~N.~Rabotnov

\book Polzuchest' elementov konstruktsiy \rm [Creep of Structural Elements]

\publaddr Moscow

\publ Nauka

\yr 1966

\totalpages 752







\RBibitem{popovkoval:bib:RadchPopov} 

\paper Построение аналитического решения стохастической краевой задачи установившейся ползучести для толстостенной трубы 

\by Н.~Н.~Попов, В.~П.~Радченко  

\jour ПММ 

\yr 2012 

\vol 76

\issue 6

\pages 1023--1031

\transl

\jour J. Appl. Math. Mech.

\by N.~N.~Popov, V.~P.~Radchenko

\paper Analytical solution of the stochastic steady-state creep boundary value problem for a thick-walled tube

\yr 2012

\vol 76

\issue 6

\pages 738--744.

\crossref{10.1016/j.jappmathmech.2013.02.011}









\RBibitem{popovkoval:bib:PopovKoval} 

\by Н.~Н.~Попов, Л.~В.~Коваленко

\paper Оценка надёжности осесимметричных стохастических элементов конструкций при ползучести по теории выбросов

\jour Вестн. Сам. гос. техн. ун-та. Сер. Физ.-мат. науки

\yr 2012

\issue 2(27)

\pages 72--77.

\mathnet{http://mi.mathnet.ru/vsgtu1090}

\crossref{10.14498/vsgtu1090}

\misctransl

\lang In Russian

\by N.~N.~Popov, L.~V.~Kovalenko

\paper Evaluation of reliability of axisymmetric stochastic elements of constructions under creepage on the basis of theory of runs

\jour Vestn. Samar. Gos. Tekhn. Univ. Ser. Fiz.-Mat. Nauki

\yr 2012

\issue 2(27)

\pages 72--77.



\RBibitem{popovkoval:bib:Sdobyrev}  

\paper Критерий длительной прочности для некоторых жаропрочных сплавов при сложном напряжённом состоянии

\by В.~П.~Сдобырев 

\jour Изв. АН СССР. ОТН. Механика и машиностроение 

\yr 1959 

\issue 6 

\pages 93--99

\misctransl

\lang In Russian

\by V.~P.~Sdobyrev

\paper The criterion of long-term strength for certain heat resisting alloys for a complex stress state

\jour Izv. AN SSSR, OTN, Mekhanika i mashinostroenie

\yr 1959 

\issue 6 

\pages 93--99.







\RBibitem{popovkoval:bib:Lokos}  

\paper Стандартизация критериев длительной прочности

\by А.~Н.~Локощенко, C.~Ф.~Шестериков

\inbook Унифицированные методы определения ползучести и длительной прочности. \rm Вып.~7

\publaddr М. 

\publ Изд-во стандартов 

\yr 1986 

\pages 3--15

\misctransl

\lang In Russian

\paper Standardization of long-term strength criteria

\by A.~N.~Lokoshchenko, C.~F.~Shesterikov

\inbook Unifitsirovannye metody opredeleniya polzuchesti i dlitel'noy prochnosti \rm [Standardized Methods for the Determination of Creep and Long-Term Strength]. Issue.~7

\publaddr Moscow

\publ Izd. Standartov

\yr 1986

\pages 3--15.











\RBibitem{popovkoval:bib:RadBash} 

\by В.~П.~Радченко, Е.~В.~Башкинова, С.~Н.~Кубышкина

\paper Об одном подходе к~оценке длительной прочности толстостенных труб на основе интегрально-средних напряженных состояний

\jour Вестн. Сам. гос. техн. ун-та. Сер. Физ.-мат. науки

\yr 2002

\issue 16

\pages 96--104.

\mathnet{http://mi.mathnet.ru/vsgtu105}

\crossref{10.14498/vsgtu105}

\misctransl

\lang In Russian

\by V.~P.~Radchenko, E.~V.~Bashkinova, S.~N.~Kubyshkina

\paper On one approach to estimate of  long-term strength for thick-walled pipe based on integral-average stress states

\jour Vestn. Samar. Gos. Tekhn. Univ. Ser. Fiz.-Mat. Nauki

\yr 2002

\issue 16

\pages 96--104.









\RBibitem{popovkoval:bib:RadchEremin} 

\by В.~П.~Радченко, Ю.~А.~Еремин 

\book Реологическое деформирование и разрушение материалов и элементов конструкций 

\publaddr М. 

\publ Машиностроение-1 

\yr 2004 

\totalpages 265

\misctransl

\lang In Russian

\by V.~P.~Radchenko, Yu.~A.~Eremin

\by Reologicheskoe deformirovanie i razrushenie materialov i elementov konstruktsiy

\rm [Rheological Strain and Fracture of Materials and Structural Elements]

\publ Mashinostroenie-1

\publaddr Moscow

\yr 2004 

\totalpages 265







\end{thebibliography}









\shortengtitle{N.~N.~P\,o\,p\,o\,v, L.~V.~K\,o\,v\,a\,l\,e\,n\,k\,o}{Reliability Evaluation of Stochastically Heterogeneous Thick-walled Pipe \dots}



\makeentitle







 \newpage

%

% \tableofengcontents

%

%\end{document}

